\chapter*{Agradecimientos}
\addcontentsline{toc}{chapter}{Agradecimientos}

El presente trabajo recoge una línea de investigación que articula control automático,
optimización y electrificación rural. Por ello, corresponde expresar un agradecimiento
especial a la Universidad Nacional del Altiplano de Puno por brindar el entorno académico
desde el cual se consolidó esta propuesta y por impulsar investigaciones vinculadas con
las necesidades energéticas reales del sur andino.

Se reconoce también el aporte de docentes, colegas y revisores que ayudaron a afinar
la formulación del problema, especialmente en la discusión sobre microredes aisladas,
turbinas Michell--Banki y criterios de desempeño para la sintonización de controladores
PID. Sus observaciones permitieron transformar una idea general en una propuesta técnica
más clara, comparable y útil para aplicaciones futuras.

Finalmente, se agradece a la población de Limbani y, en general, a las comunidades
rurales que inspiran este tipo de estudios. La búsqueda de soluciones de control
robustas, económicas y adaptables solo adquiere sentido cuando se vincula con la
continuidad del servicio eléctrico, la protección de los equipos y la sostenibilidad de
los sistemas que sostienen la vida cotidiana.
